\documentclass[runningheads]{llncs}

\usepackage{graphicx}
\usepackage{amsmath}
\usepackage{amssymb}
\usepackage{booktabs}
\usepackage{tikz}
\usepackage{pgfplots}
\pgfplotsset{compat=1.17}
\usepackage{xcolor}
\usepackage{listings}
\usepackage{hyperref}

% llncs does not define \doi, but splncs04.bst emits it.
\providecommand{\doi}[1]{\textsc{doi}:\ \texttt{#1}}

\definecolor{codebg}{rgb}{0.96,0.96,0.96}
\lstdefinelanguage{Coq}{
  morekeywords={Lemma,Theorem,Proof,Qed,Defined,forall,intros,rewrite,apply,
                reflexivity,symmetry,induction,simpl,lia,Fail,Abort,
                Definition,Fixpoint,Inductive,Class,Instance,trocq,Trocq,Use},
  sensitive=true,
  morecomment=[s]{(*}{*)},
  commentstyle=\color{gray}\itshape,
  keywordstyle=\color{blue}\bfseries,
  basicstyle=\ttfamily\small,
  breaklines=true,
}
\lstset{language=Coq,backgroundcolor=\color{codebg},frame=single,columns=fullflexible}

\begin{document}

\title{When Does Proof Transfer Pay Off?\\
An Empirical Cost Model for Trocq versus Manual\\
Rewriting in Coq}

\author{Charles Santos\inst{1} \and Vander Alves\inst{1} \and Leopoldo Teixeira\inst{2}}
\authorrunning{C. Santos, V. Alves, L. Teixeira}

\institute{Department of Computer Science, University of Bras\'{i}lia (UnB),
Bras\'{i}lia, Brazil\\
\email{thyarles@gmail.com}
\and
Informatics Center, Federal University of Pernambuco (UFPE), Recife, Brazil}

\maketitle

\begin{abstract}
Interactive theorem provers routinely maintain the same mathematical concept
under several representations---for instance, a proof-oriented monomorphic
data structure alongside a computation-oriented polymorphic one---and proving
theorems on one representation obliges a proof engineer to either replay the
argument or \emph{transfer} it from the other. Proof-transfer frameworks
built on (univalent) parametricity, such as Trocq, promise to automate this
replay, but automation is not free: it trades a per-theorem cost for a fixed
setup cost, and it is not obvious a priori when that trade is worthwhile.
This paper reports a small, fully machine-measured case study contrasting
manual proof transfer against Trocq-based transfer on monomorphic/polymorphic
data structures in Coq/Rocq, and develops an economic (return-on-investment)
model of the trade-off through four successive refinements, each falsified and
corrected by evidence read mechanically off the Coq scripts. The decisive
refinement is the fourth: against the baseline a working engineer actually
uses---copy-paste with find-and-replace---Trocq does not break even until
roughly fifteen transferred theorems, nearly five times later than the
bridge-based baseline suggests. We show that this copy-paste baseline is
available only while the two representations share an algebraic signature,
and that it collapses as soon as the signature diverges, at which point
mechanical replication stops being mechanical. The resulting model is
therefore not a single break-even curve but a \emph{regime classification},
and we argue the classification, rather than the fitted constants, is what
transfers to larger developments. All constants are regenerated from the
artifact by a counting script, so every number in the paper is traceable to a
line of Coq.

\keywords{Proof transfer \and Trocq \and univalent parametricity \and Coq/Rocq \and
proof engineering cost \and empirical software engineering}
\end{abstract}


\section{Introduction}
\label{sec:intro}

Formal developments in interactive theorem proving frequently keep more than
one representation of the same underlying concept alive at once. A classic
example is a pair of list types: one monomorphic and convenient to reason
about, the other polymorphic (or otherwise more general) and convenient to
compute with or reuse across a library. Once both representations exist,
every theorem proved about one of them is, semantically, also true of the
other---yet formally, nothing forces the two proof developments to stay in
sync. Restating and reproving each theorem by hand is tedious and does not
scale; this is the classical \emph{proof-transfer} problem, and it recurs
throughout formalized mathematics and program verification, including in
software-product-line (SPL) settings~\cite{apel2013featureoriented,thum2014classification}
where the same variability-aware property must often be established over both
a feature-oriented and a product-level view of a system.

Proof assistants have accumulated a family of tools that attack this
problem automatically, ranging from data refinement frameworks such as
CoqEAL~\cite{cohen2013refinements} and the Isabelle/HOL Lifting and Transfer
package~\cite{huffman2013lifting}, to frameworks grounded in
\emph{univalent parametricity}~\cite{tabareau2018equivalences,tabareau2021marriage}.
Trocq~\cite{cohen2024trocq,cohen2025trocq} is the most recent and most
general of the latter family: it generalizes the univalent parametricity
translation to a hierarchy of relatedness classes so that a single
\texttt{trocq} tactic call can, in favorable cases, discharge most of a
proof obligation automatically, leaving only a residual lemma about the
target representation.

Automation of this kind is attractive, but it is not obviously
\emph{always} the right engineering choice. Setting up a Trocq translation
for a given pair of types requires writing relational witnesses and
registering them with the plugin---work that must be paid regardless of how
many theorems will eventually benefit from it. Manual transfer, by
contrast, avoids that setup but repeats some argument for every theorem.
Whether Trocq ``pays off'' therefore depends on how much proof-transfer work
is actually being done, and on how the two approaches' costs scale with the
size and shape of that work. The proof-engineering literature has begun to
treat such questions empirically---the QED-at-Large survey~\cite{ringer2020qed}
catalogues effort and maintenance concerns across verified-software
projects---but, to the best of our knowledge, no explicit cost model has been
published for proof transfer specifically, where the literature has so far
focused on expressiveness and correctness rather than on effort.

Our core scientific hypothesis is that the initial setup overhead required by
univalent parametricity tools like Trocq is offset by a lower per-theorem
cost compared to manual rewriting, and that the break-even point is governed
by the structural complexity of the transferred theorems and the size of the
vocabulary shared between the representations. As will become clear, the
study falsifies the second half of this hypothesis and replaces it with a
sharper one.

This paper makes the following contributions:
\begin{itemize}
\item A small, controlled Coq/Rocq case study contrasting manual proof
      transfer against Trocq-based transfer, instrumented with a syntactic
      effort proxy---tactic steps---that is \emph{mechanically} extracted by
      a counting script rather than tallied by hand
      (Section~\ref{sec:setup}).
\item An economic cost model of the manual-versus-Trocq trade-off,
      developed through four iterations, each refuted and corrected by
      evidence read from the scripts rather than assumed a priori
      (Section~\ref{sec:model}). The iterative structure is itself a
      methodological contribution: it shows how a naive cost intuition about
      a proof-automation tool can be wrong in ways that only mechanical
      inspection reveals.
\item The central empirical finding: the manual baseline is not one baseline
      but two, and which one applies is decided by whether the two
      representations share an algebraic signature
      (Section~\ref{sec:fourth}). We give the break-even point in both
      regimes and characterize the boundary between them.
\item A decision heuristic organized around this regime classification, and a
      scalability argument defending the model's applicability beyond the
      running example (Section~\ref{sec:results}).
\item A threats-to-validity discussion following an \emph{evidence-debt}
      framing~\cite{alves2025evidencedebt}, treating unresolved concerns as
      explicit, discharge-able obligations (Section~\ref{sec:ttv}).
\end{itemize}

We stress at the outset that the case study is deliberately small. That is a
design choice, not only a limitation: a corpus small enough to be measured
exhaustively and re-measured mechanically is what allows every constant below
to be traced to a specific line of Coq. Section~\ref{sec:scope} argues
explicitly for why the resulting classification, as opposed to the fitted
constants, is what we expect to carry over to larger developments.


\section{Background}
\label{sec:background}

\subsection{Proof transfer and its cost}
Given two types $A$ and $B$ known to be related---typically by an
isomorphism, or more generally by some relation $R$---proof transfer aims
to derive a theorem about $B$ from an already-proved theorem about $A$
(or vice versa) without re-deriving it from scratch. The problem is folklore
in mathematics, where such transfers are usually left implicit
(``clearly, the same argument applies''), but in a proof assistant every
step must be justified, and the bridging argument becomes explicit proof
work. That work has two natural components: a \emph{setup} cost, incurred
once per pair of representations and per piece of shared vocabulary
(functions, constructors) that both sides use; and a \emph{per-theorem}
cost, incurred once for every theorem that is actually transferred.

\subsection{Univalent parametricity and Trocq}
Parametricity~\cite{reynolds1983types} associates to a type a relation
describing how its elements are related across two interpretations of the
same type former; \emph{univalent} parametricity~\cite{tabareau2018equivalences,tabareau2021marriage}
strengthens ordinary parametricity so that it is preserved by type
equivalences---appealing, in its original formulation, to the univalence
axiom~\cite{hottbook}---which lets a proof engineer transport definitions and
theorems along an equivalence between $A$ and $B$ largely automatically.
Extending parametricity to inductive types and to the propositional fragment
raises subtleties of its own~\cite{anand2017revisiting}.
Trocq~\cite{cohen2024trocq,cohen2025trocq} generalizes this translation
into a calculus with several \emph{relatedness classes} (Param44 for
isomorphisms down to weaker classes for non-invertible relations), together
with a Coq/Rocq plugin, implemented in Coq-Elpi~\cite{tassi2019elpi}, that
exposes a \texttt{trocq} tactic and \texttt{Trocq Use} vernacular commands
for registering the relational witnesses of a given translation. Once a
translation is registered, \texttt{trocq} rewrites a goal about one
representation into a goal about the other, discharging as much of the
argument as the calculus can justify automatically.

One implementation detail matters for the cost model, because it turns out to
be a recurring and measurable overhead. Trocq's witness database is keyed on
the head \emph{global reference} of a term. A relational witness whose
conclusion is headed by a locally bound variable---as happens whenever the
witness is stated polymorphically over the element type---cannot be
registered, and must be duplicated at each concrete instantiation. We
quantify the resulting cost in Section~\ref{sec:scope}.

\subsection{Type classes and algebraic signatures}
Coq's type classes~\cite{sozeau2008first} let a polymorphic definition
require algebraic structure of its element type: our polymorphic summation
takes an \texttt{Addable} instance supplying a binary operation, a unit, and
their laws. Whether two representations agree on such a signature turns out
to decide which cost regime a development is in
(Section~\ref{sec:fourth}), so we make the notion explicit here: two
representations \emph{share an algebraic signature} when every operation a
theorem mentions is available on both sides with the same laws, so that
renaming one side's proof script yields a script that still typechecks.

\subsection{Related transfer and repair mechanisms}
CoqEAL's \emph{Refinements for free}~\cite{cohen2013refinements} pursues a
related but distinct goal: transporting \emph{computational} content
(e.g., efficient data representations) rather than proofs, using
parametricity translations tailored to refinement relations. The
Isabelle/HOL Lifting and Transfer package~\cite{huffman2013lifting}
addresses the analogous problem for quotient types in a different proof
assistant. \emph{Proof repair}~\cite{ringer2018proof,ringer2021proofrepair,viola2025proof}
takes a complementary angle: rather than transferring a proof to a new
representation ahead of time, it patches an existing proof after a
representation change breaks it, using type-equivalence or quotient-type
equivalence information much like Trocq's relational witnesses; related
automation repairs verified libraries after
refactoring~\cite{gopinathan2023mostly}. None of these lines of work, to our
knowledge, reports an explicit cost model comparing automated transfer
against the manual alternative it is meant to replace.

\subsection{Cost modeling}
Software cost estimation has a long history of decomposing effort into a
size-independent (fixed) component and a size-dependent (variable)
component, most influentially in Boehm's COCOMO
model~\cite{boehm1981economics}, which fits such a decomposition to
project size and a set of scaling factors. We borrow this fixed/variable
decomposition, applying it not to lines of code but to proof obligations,
and use it to structure the manual-versus-Trocq comparison in
Section~\ref{sec:model}.


\section{Case Study Setup}
\label{sec:setup}

\subsection{The corpus}

The case study formalizes several pairs of representations in Coq/Rocq
(Table~\ref{tab:files}). The core pair is a monomorphic list of natural
numbers, \texttt{NatList}, with operations \texttt{nlength} and \texttt{napp}
(and later \texttt{nrev}, \texttt{nsum}), alongside a polymorphic list
\texttt{PList A} with \texttt{plength}, \texttt{papp}, \texttt{prev} and a
type-class-mediated \texttt{psum}. Two further pairs probe the model outside
lists: an arithmetic expression language (\texttt{bs\_z1}, \texttt{bs\_z1\_lift})
and a higher-order, function-encoded map type (\texttt{bs\_m1}), whose map
interface---\texttt{total\_map}, \texttt{tm\_update}, \texttt{partial\_map}
and the \texttt{tm\_update\_neq} property---follows the treatment in
\emph{Software Foundations}~\cite{softwarefoundations}. A separate experiment
(\texttt{bs\_f1}) transfers a Fermat-style statement from $\mathbb{Z}$ to
$\mathbb{Z}/9\mathbb{Z}$, where the purpose of transfer is decidability
rather than effort saving.

The corpus comprises 15 files, 194 machine-checked proofs, 884 tactic steps
and 85 \texttt{Trocq Use} registrations. All files but one compile under the
project build; the exception is discussed in Section~\ref{sec:ttv}.

\begin{table}[t]
\centering
\caption{Roles of the Coq/Rocq scripts making up the case study. ``Steps''
are tactic steps as measured by the counting script of
Section~\ref{sec:convention}.}
\label{tab:files}
\begin{tabular}{@{}llr@{}}
\toprule
File & Role & Steps \\
\midrule
\multicolumn{3}{@{}l}{\emph{Baby Step 1 --- monomorphic versus polymorphic lists}}\\
\texttt{bs\_p1.v} & \texttt{NatList}: monomorphic source, base theorems & 19 \\
\texttt{bs\_p2.v} & \texttt{PList A}: polymorphic target, same theorems & 9 \\
\texttt{bs\_p3.v} & The same proofs as raw proof terms (recursors) & 0 \\
\texttt{bs\_p4.v} & Manual transfer: explicit bijections and bridges & 26 \\
\texttt{bs\_p5.v} & Trocq transfer: relational witnesses, \texttt{trocq} & 44 \\
\texttt{bs\_p6.v} & Adds \texttt{rev}: the complexity benchmark & 46 \\
\texttt{bs\_p7.v} & Adds \texttt{sum}: the simplicity counterpoint & 28 \\
\midrule
\multicolumn{3}{@{}l}{\emph{Baby Step 2 --- copy-paste baseline and signature divergence}}\\
\texttt{bs\_a1.v} & ROI experiment against a \emph{copy-paste} baseline & 138 \\
\texttt{bs\_a2.v} & \texttt{Addable} type class: where copy-paste breaks & 200 \\
\texttt{bs\_a3.v} & Parametric infrastructure; fairness of the comparison & 150 \\
\texttt{bs\_a4.v} & Three strategies isolated on \texttt{length} alone & 62 \\
\texttt{bs\_z1.v} & Minimized arithmetic base language & 9 \\
\texttt{bs\_z1\_lift.v} & Divergence test: a copy-pasted proof that fails & 59 \\
\texttt{bs\_m1.v} & Higher-order maps: scalability probe & 52 \\
\texttt{bs\_f1.v} & $\mathbb{Z}\to\mathbb{Z}/9\mathbb{Z}$: space reduction & 42 \\
\bottomrule
\end{tabular}
\end{table}

\subsection{The effort proxy and its counting convention}
\label{sec:convention}

As an effort proxy we count \emph{proof obligations}, operationalized as
tactic steps, under a convention fixed once and applied mechanically:

\begin{itemize}
\item one period-terminated tactic counts as one step;
\item a chained tactic \texttt{t1; t2.} counts as one step;
\item bullet markers (\texttt{-}, \texttt{+}, \texttt{*}, braces) count as zero;
\item \texttt{Trocq Use} commands are counted separately, as commands, and
      are charged to Trocq's setup at one step each.
\end{itemize}

Crucially, the counts are not tallied by hand. A script,
\texttt{\_count\_steps.py}, parses every file, strips nested comments and
string literals, isolates proof bodies, and emits a machine-readable table
from which every constant in this paper is drawn. This matters more than it
may appear: an earlier, hand-tallied revision of this study disagreed with
the mechanical count in three files, once by reporting a nine-step proof that
is in fact eleven steps, and once by reporting a shared subtotal of 20 where
the measured value is 33. Those discrepancies were not visible to inspection;
they were visible only to the parser.

This is a coarse, syntactic measure. It does not weight tactics by search
cost or by the cognitive effort of discovering them, and
Section~\ref{sec:ttv} returns to this limitation with a concrete instance of
it drawn from the corpus itself. Its virtue is that it is unambiguous and
reproducible: every parameter below can be regenerated from the artifact by
re-running one script.

\subsection{Guiding question}

As more theorems are transferred, and as those theorems grow more complex, at
what point---if any---does Trocq's per-theorem economy overcome its setup
cost, and what determines where that point lies?


\section{Iterative Development of the Cost Model}
\label{sec:model}

We present the cost model as four iterations. Each is stated, then checked
against the corpus; each failure motivates the next. We keep this structure,
rather than presenting only the final model, because the failures are
informative in their own right: they show three specific ways in which a
plausible intuition about proof-automation cost is wrong, and the third is
the paper's main finding.

\subsection{First iteration: a purely linear model}
\label{sec:first}

The simplest model assumes both approaches' costs scale linearly with the
number of theorems transferred, $n$, and with a per-theorem complexity factor
$c$, and assumes Trocq alone pays a one-time setup cost governed by a
type-distance factor $d$ (with $d=1$ for the isomorphism class \emph{Param44}
used throughout):
\begin{align}
C_{\text{manual}}(n,c) &= 7cn \\
C_{\text{trocq}}(n,c,d) &= (12d+5) + 2cn
\end{align}
With $c=1$, $d=1$ these give a break-even at $n^{*}=17/5\approx 3.4$
theorems.

\paragraph{Why this is wrong.} The model assumes manual transfer has
\emph{no} fixed cost. Inspecting \texttt{bs\_p4.v} shows this is false: the
bijections between \texttt{NatList} and \texttt{PList}, and the ``bridge''
lemmas relating \texttt{plength}/\texttt{papp} to their monomorphic
counterparts, are written once, not once per theorem.

\begin{lstlisting}
(* bs_p4.v -- written once, not per theorem *)
Lemma plist_natlist_iso : natlist_to_plist (plist_to_natlist l) = l.
Lemma natlist_plist_iso : plist_to_natlist (natlist_to_plist l) = l.
Lemma plength_eq_nlength : nlength (plist_to_natlist l) = plength l.
Lemma plist_to_natlist_app :
  plist_to_natlist (papp l1 l2) = napp (plist_to_natlist l1) (plist_to_natlist l2).
\end{lstlisting}

Moreover \texttt{bs\_p5.v} shows Trocq's own base cost is not a pure fixed
intercept either: it requires \emph{the same} bridge lemmas, restated for the
relational encoding. Both halves of the split were mischaracterized.

\subsection{Second iteration: separating shared setup from per-theorem cost}
\label{sec:second}

The second iteration factors out a \emph{shared} base cost that both
approaches pay and that scales with $f$, the number of distinct functions the
theorems mention (here $f=2$: \texttt{plength} and \texttt{papp}):
\begin{align}
C_{\text{base}}(f) &= S_{\text{iso}} + f \cdot S_{\text{bridge}} \\
C_{\text{manual}}(n,f) &= C_{\text{base}}(f) + n \cdot P_{\text{manual}} \\
C_{\text{trocq}}(n,f) &= C_{\text{base}}(f) + f \cdot W_{\text{trocq}} + n \cdot 2
\end{align}
Measured on \texttt{bs\_p4.v}/\texttt{bs\_p5.v}: the mutual-inverse proofs
cost $S_{\text{iso}}=10$ steps; the two bridge lemmas cost $5$ and $6$, so
$f\cdot S_{\text{bridge}}=11$ and $C_{\text{base}}(2)=21$. Trocq's
\emph{additional} overhead is the \texttt{Param44} packaging
(\texttt{R\_NatList}, $5$ steps), three shared registrations, and one
relational witness plus one registration per function
($5{+}1$ for \texttt{plength}, $7{+}1$ for \texttt{papp}), totalling $22$,
i.e.\ $W_{\text{trocq}}\approx 11$ per function. The per-theorem manual
rewrite chain in \texttt{bs\_p4.v} measures $P_{\text{manual}}=7$, against a
flat two-tactic close in \texttt{bs\_p5.v}:

\begin{lstlisting}
(* bs_p5.v -- per-theorem cost is just 2 tactics *)
Theorem _plength_papp : forall (l1 l2 : _PList),
  _plength (_papp l1 l2) = _plength l1 + _plength l2.
Proof.
  trocq.
  apply nlength_napp.
Qed.
\end{lstlisting}

This gives $C_{\text{manual}}(n) = 21+7n$ and $C_{\text{trocq}}(n) = 43+2n$:
Trocq starts more expensive but overtakes manual transfer at
$n^{*}=22/5=4.4$ theorems. Its return on investment converges, as
$n\to\infty$, to the finite constant
\begin{equation}
\mathrm{ROI}_{\infty} = \frac{P_{\text{manual}} - 2}{2},
\end{equation}
which for $P_{\text{manual}}=7$ gives a bounded $2.5\times$ return---not an
ever-growing one, because both approaches' \emph{marginal} costs are
constant; only the intercepts differ.

\subsection{Third iteration: complexity-sensitive theorems}
\label{sec:third}

The second iteration still assumes every theorem costs the same
$P_{\text{manual}}$ to transfer manually. \texttt{bs\_p6.v}, which adds
\texttt{rev}, falsifies this. The theorem \texttt{\_prev\_papp} (reversing a
concatenation) mentions five function applications (\texttt{\_prev} three
times, \texttt{\_papp} twice) and costs \emph{eleven} manual tactic steps,
while its Trocq proof remains exactly two:

\begin{lstlisting}
(* Manual -- 11 tactic steps *)
Theorem _prev_papp_manual :
  _prev (_papp l1 l2) = _papp (_prev l2) (_prev l1).
Proof.
  intros l1 l2.                                                (*  1 *)
  rewrite <- (plist_nlist_iso (_prev (_papp l1 l2))).          (*  2 *)
  rewrite <- (plist_nlist_iso (_papp (_prev l2) (_prev l1))).  (*  3 *)
  apply f_equal.                                               (*  4 *)
  rewrite plist_2_nlist_rev.                                   (*  5 *)
  rewrite plist_2_nlist_app.                                   (*  6 *)
  rewrite nrev_napp.                                           (*  7 *)
  rewrite plist_2_nlist_app.                                   (*  8 *)
  rewrite plist_2_nlist_rev.                                   (*  9 *)
  rewrite plist_2_nlist_rev.                                   (* 10 *)
  reflexivity.                                                 (* 11 *)
Defined.

(* Trocq -- 2 tactic steps *)
Theorem _prev_papp_trocq :
  _prev (_papp l1 l2) = _papp (_prev l2) (_prev l1).
Proof.
  trocq.
  apply nrev_napp.
Defined.
\end{lstlisting}

We therefore introduce $c_{\text{avg}}$, the average number of function
applications per transferred theorem, and a constant $k$ relating complexity
to manual tactic count; fitting $11$ steps over $5$ applications gives
$k=2.2$. Trocq's per-theorem cost, crucially, does \emph{not} grow with
$c_{\text{avg}}$ in this data: the \texttt{trocq} tactic dispatches the
relational bookkeeping regardless of how many functions the goal mentions:
\begin{align}
C_{\text{manual}}(n,f,c_{\text{avg}}) &= C_{\text{base}}(f) + n \cdot (k \cdot c_{\text{avg}}) \\
C_{\text{trocq}}(n,f) &= C_{\text{base}}(f) + f \cdot W_{\text{trocq}} + n \cdot 2 \\
\mathrm{ROI}(n,f,c_{\text{avg}}) &= \frac{n(k \cdot c_{\text{avg}} - 2) - f \cdot W_{\text{trocq}}}{C_{\text{base}}(f) + f \cdot W_{\text{trocq}} + 2n}
\end{align}
With \texttt{rev} included, $f=3$, $S_{\text{iso}}=10$ and the three bridges
cost $5+6+6=17$, so $C_{\text{base}}(3)=27$; Trocq's extra is
$5+3+6+8+6=28$. Hence $C_{\text{manual}}(n) = 27+11n$ and
$C_{\text{trocq}}(n) = 55+2n$, with a break-even at
$n^{*} = 28/9 \approx 3.1$ theorems, down from $4.4$. Complexity lowers the
break-even; it does not abolish it. The asymptotic return rises to
$(k\cdot c_{\text{avg}}-2)/2 = 4.5$.

\paragraph{Why this is still wrong.} Every iteration so far has compared
Trocq against a manual transfer that routes through explicit bijections. But
that is not what a working proof engineer does first.

\subsection{Fourth iteration: the copy-paste baseline, and where it collapses}
\label{sec:fourth}

Confronted with a theorem proved for \texttt{NatList} and needed for
\texttt{PList nat}, the first thing an engineer tries is not a bijection. It
is to copy the proof script and rename the identifiers. \texttt{bs\_a1.v}
takes this adversarial baseline seriously and measures it.

The catch is that the baseline's viability is a property of how the target
was defined. In \texttt{bs\_a1.v} the polymorphic operations are deliberately
declared at \texttt{PList nat}---monomorphic---so that every proof is a
literal rename of its \texttt{NatList} counterpart:

\begin{lstlisting}
(* bs_a1.v -- Phase 2: NatList proof, renamed. Nothing else changes. *)
Theorem plength_papp_manual : forall (l1 l2 : PList nat),
    plength (papp l1 l2) = plength l1 + plength l2.
Proof.
    intros l1 l2.
    induction l1 as [| h t IH]; simpl.
    - reflexivity.
    - rewrite IH. reflexivity.
Defined.
\end{lstlisting}

Under this baseline the manual side has \emph{no} fixed cost at all: no
bijections, no bridges, no glue. Measured over three theorems ($5$, $5$ and
$7$ steps) the average is $P_{\text{paste}} = 17/3 \approx 5.67$, and:
\begin{align}
C_{\text{manual}}(n) &= P_{\text{paste}}\cdot n \;\approx\; 5.67\,n \\
C_{\text{trocq}}(n)  &= 53 + 2n
\end{align}
where Trocq's $53$ decomposes, measured, as $16$ for the isomorphism block
(two mutual-inverse proofs, the \texttt{Param44} packaging and three shared
registrations) plus $11+14+12$ for the three functions. The break-even moves
to
\begin{equation}
n^{*} = \frac{53}{17/3 - 2} = \frac{159}{11} \approx 14.5 ,
\end{equation}
and the asymptotic return collapses to
$\mathrm{ROI}_{\infty} = (17/3-2)/2 = 11/6 \approx 1.83$. Against the
baseline an engineer actually reaches for, Trocq must amortize its setup over
roughly \emph{fifteen} theorems---almost five times the third iteration's estimate
and an order of magnitude worse than the first iteration's.

\paragraph{The boundary.} This would be a damning result for proof transfer
if the copy-paste baseline were always available. It is not. Renaming
succeeds precisely while the two representations share an algebraic
signature. \texttt{bs\_a2.v} removes that assumption: \texttt{PList} becomes
genuinely polymorphic and summation is mediated by a type class,
\begin{lstlisting}
Class Addable (A : Type) : Type := {
    add : A -> A -> A;   zero : A;
    add_assoc  : forall x y z : A, add (add x y) z = add x (add y z);
    add_zero_l : forall x : A, add zero x = x;
    add_zero_r : forall x : A, add x zero = x }.
\end{lstlisting}
and the renamed proof no longer typechecks. Two steps must be repaired by
hand, each with reasoning specific to the signature rather than to the
theorem:
\begin{lstlisting}
Theorem psum_papp_manual :
    forall {A : Type} {H : Addable A} (l1 l2 : PList A),
    psum (papp l1 l2) = add (psum l1) (psum l2).
Proof.
    intros A H l1 l2.
    induction l1 as [| h t IH]; simpl.
    - symmetry. apply add_zero_l.        (* DIVERGENCE 1: was [reflexivity] *)
    - rewrite IH. symmetry. apply add_assoc.  (* DIVERGENCE 2: was [lia]    *)
Defined.
\end{lstlisting}

\texttt{bs\_z1\_lift.v} sharpens this from repair to outright failure. There
the \texttt{Addable} signature provides only a right identity while
evaluation associates right-to-left, so the copied goal reduces to one the
class cannot discharge. The file records the failure as a checked artifact:

\begin{lstlisting}
(* Addable provides only add_zero_r; the goal needs left identity. *)
Goal forall {A : Type} {H : Addable A} (e : pexpr A),
  peval (PPlus e (PConst zero)) = peval e.
Proof. intros. simpl. Fail apply add_zero_r. Abort.
\end{lstlisting}

No rename rescues this; the engineer must supply new mathematics. Trocq, by
contrast, absorbs the same gap once, inside a single bridge lemma whose
inductive step closes with \texttt{lia}, after which all three theorems in
that file transfer in two steps each.

\paragraph{The model is a regime classification.} The fourth iteration's
lesson is not a different constant but a different shape. There is no single
break-even curve, because there is no single manual baseline
(Table~\ref{tab:regimes}). What decides which curve applies is whether
renaming still produces a script that typechecks.

\begin{table}[t]
\centering
\caption{The two regimes. All constants measured by
\texttt{\_count\_steps.py} over the corpus.}
\label{tab:regimes}
\begin{tabular}{@{}lll@{}}
\toprule
 & \textbf{Regime A} & \textbf{Regime B} \\
\midrule
Condition & signature shared & signature diverges \\
Manual method & copy-paste and rename & bijections and bridges \\
Evidence & \texttt{bs\_a1.v} & \texttt{bs\_p4.v}, \texttt{bs\_a2.v}, \texttt{bs\_z1\_lift.v} \\
$C_{\text{manual}}(n)$ & $5.67\,n$ & $27 + 11n$ \\
$C_{\text{trocq}}(n)$ & $53 + 2n$ & $55 + 2n$ \\
Break-even $n^{*}$ & $\approx 14.5$ & $\approx 3.1$ \\
$\mathrm{ROI}_{\infty}$ & $\approx 1.83\times$ & $\approx 4.5\times$ \\
Verdict & Trocq rarely worth it & Trocq worth it early \\
\bottomrule
\end{tabular}
\end{table}


\section{Results and a Decision Heuristic}
\label{sec:results}

\subsection{The heuristic}
\label{sec:heuristic}

The model yields a decision procedure for a proof engineer choosing, before
writing anything, whether to invest in a Trocq translation for a given pair
of representations.

\begin{enumerate}
\item \textbf{Determine the regime first.} Attempt the rename on the single
      hardest theorem you expect to transfer. If it typechecks, you are in
      Regime A and your baseline is nearly free; budget for $n^{*}\approx 15$
      before Trocq repays its setup. If it fails---because an operation is
      mediated by a type class, because the two sides' laws differ, or
      because evaluation order or associativity diverges---you are in Regime
      B, where the baseline itself costs bridges and Trocq repays within
      three or four theorems. This test is cheap, mechanical, and dominates
      every other consideration below.
\item \textbf{Then estimate the shared vocabulary $f$.} Trocq's setup is
      paid per function, not per theorem, so a large vocabulary relative to
      the number of theorems actually needed argues against it. In Regime A
      this compounds with an already-late break-even.
\item \textbf{Then weigh complexity.} If theorems are structurally complex
      (many nested applications, as with \texttt{rev} over \texttt{app}),
      the manual per-theorem cost rises as $k\cdot c_{\text{avg}}$ while
      Trocq's stays at two, pulling $n^{*}$ down. Complexity moves the
      break-even; it does not, on this data, eliminate it.
\item \textbf{Do not count on cross-type amortization.} We had expected
      Trocq's setup to amortize across element types. It does not, in this
      corpus: see Section~\ref{sec:scope}.
\end{enumerate}

\subsection{Scope, scalability, and what generalizes}
\label{sec:scope}

The running example is a list of natural numbers. It is fair to ask whether
anything established on it survives contact with a realistic development. We
take the objection seriously and answer it in four parts, conceding the first.

\paragraph{The example is simple, deliberately.} A corpus small enough to
measure exhaustively is what makes the constants auditable: every number in
this paper is regenerated by one script from 194 proofs, and the discrepancies
that script found in our own earlier hand-tallies
(Section~\ref{sec:convention}) are precisely the errors a larger, unmeasurable
corpus would have concealed. We would rather report a small result that is
mechanically checkable than a large one that is not. What we claim
generalizes is not the fitted constants but the regime classification, and
the mechanism behind it.

\paragraph{The mechanism scales with signature divergence, not with size.}
The quantity that decides the regime is not how big the development is but
how many \emph{diverging algebraic obligations} separate the two
representations. Manual replication must discharge each divergence
separately, and must do so again for every theorem that touches it: the cost
grows with the product of theorems and divergences. Trocq's setup, in
contrast, is paid once per function regardless of how many theorems follow---%
each divergence is absorbed into a single relational witness. In
\texttt{bs\_a2.v} the \texttt{Addable} signature contributes two divergences
in one theorem; in a development where the source is a proof-oriented
structure and the target carries a full algebraic hierarchy, the divergence
count is a property of the hierarchy, and the manual side's cost scales with
it multiplicatively while Trocq's does not. This is the asymmetry we expect
to persist, and it is the reason a simple example can license a claim about
larger ones: the example fixes the mechanism, and the mechanism is
size-independent.

\paragraph{The counterweight, measured.} Honesty requires reporting the
mechanism that pushes the other way, and our corpus quantifies two instances.
First, Trocq's database is keyed on head global references, so a relational
witness stated polymorphically cannot be registered. In \texttt{bs\_a2.v},
four such generic witnesses are written, correctly, at a cost of $22$ tactic
steps---and then cannot be used, forcing a duplicate set of concrete
witnesses. That is $22$ of the file's $104$ setup steps spent on nothing but
the tool's addressing scheme. Second, the cost of setup grows with the
complexity of the type former, not only with the number of functions:
\texttt{bs\_m1.v}, which transfers over function-encoded maps
(\texttt{total\_map A = string -> A}), requires \emph{35} \texttt{Trocq Use}
registrations against \texttt{bs\_p5.v}'s five, because the relatedness
lattice must be registered by hand for each of five participating types. A
development rich in higher-order type formers will pay this repeatedly. Both
effects raise $n^{*}$, and neither is visible in the list example alone.

\paragraph{Cross-type amortization does not hold.} We expected the setup to
amortize when a second element type was added, and it does not. In
\texttt{bs\_a2.v}, transferring the same theorems to \texttt{PList Z} costs
Trocq a fresh set of concrete witnesses and registrations ($24$ steps, or $30$
including the transferred theorems themselves)---whereas a manual proof that was \emph{stated polymorphically} discharges the
same theorems in $9$ steps in total. Two of those
theorems (\texttt{papp\_assoc\_Z}, \texttt{prev\_papp\_Z}) are literally one
\texttt{apply}. The lesson is that Trocq's advantage in this corpus is not
reuse across element types; it is the absorption of signature divergence.
Where a manual proof can be stated polymorphically, it generalizes for free
and Trocq should not be expected to compete.

\paragraph{Beyond lists.} Two files test whether the machinery is
list-shaped. \texttt{bs\_z1\_lift.v} runs the whole argument on an arithmetic
expression language, where the divergence is in evaluation order rather than
in a data structure, and reproduces the Regime B pattern. \texttt{bs\_f1.v}
transfers a Fermat-style statement from $\mathbb{Z}$ to
$\mathbb{Z}/9\mathbb{Z}$, collapsing an infinite search space to $729$ cases;
there the value of transfer is decidability, and the cost model, which prices
effort, simply does not apply---a boundary we state rather than paper over.


\section{Threats to Validity}
\label{sec:ttv}

Following an evidence-debt framing~\cite{alves2025evidencedebt}, we state each
threat as an open obligation: a checkable condition under which a claim would
need revision, together with what would discharge it.

\paragraph{Construct validity (O1): tactic count as an effort proxy.}
That tactic steps approximate proof \emph{effort} is a modeling choice. It
ignores tactic search cost, the effort of discovering a relational witness,
and the cost of debugging a failed \texttt{trocq} call. Our corpus contains a
sharp demonstration of the proxy's blindness: \texttt{bs\_p3.v} proves the
same theorems as \texttt{bs\_p1.v} and \texttt{bs\_p2.v} but writes them as
raw proof terms applied to recursors, with no \texttt{Proof} blocks at all.
It therefore measures \emph{zero} steps while representing real work. A
second instance is visible in \texttt{bs\_p7.v}, where the same theorem is
proved in seven steps and in five, the difference being fluency with a
repetition combinator rather than any change in difficulty. Discharging this
obligation requires re-running the study with an independent proxy---proof-term
size, or time-to-completion under a controlled protocol---and checking that
the regime classification of Table~\ref{tab:regimes} is preserved.

\paragraph{Internal validity (O2): single-author, single-domain parameters.}
The fitted constants come from scripts written by one author in one domain.
An author fluent in Trocq's idioms may write shorter witnesses than a
newcomer, biasing $W_{\text{trocq}}$ downward; the seven-versus-five variance
noted above is a measured instance of exactly this effect. Discharge requires
a multi-author replication with scripts written independently and blind to
the cost model.

\paragraph{External validity (O3): relatedness class and domain.}
The study exercises the \emph{Param44} (isomorphism) class on lists,
expressions and maps. Trocq extends to weaker, non-invertible relations,
where the trade-off may differ qualitatively---$W_{\text{trocq}}$ may grow
with the relational distance $d$ in ways the first iteration's discarded
$d$-term gestured at. Discharge requires repeating
Sections~\ref{sec:model}--\ref{sec:results} on at least one lower relatedness
class and one substantially larger structure, ideally a variability-aware
artifact drawn from a software product line~\cite{thum2014classification},
and checking that $n^{*}$ and $\mathrm{ROI}_{\infty}$ retain their functional
form.

\paragraph{Reliability (O4): tool-version dependence.} The counts depend on
the Trocq plugin version, whose automation is under active development. A
later version could shrink $W_{\text{trocq}}$ or the two-tactic residual,
moving every break-even in Trocq's favor without any change to the manual
baseline. In particular, if the gref-keying limitation quantified in
Section~\ref{sec:scope} were lifted, $22$ of \texttt{bs\_a2.v}'s $104$ setup
steps would disappear outright. Discharge requires pinning and reporting
versions in the replication package and re-measuring on each release.

\paragraph{Completeness (O5): one unverified sub-claim.}
\texttt{bs\_a3.v} was written to provide a like-for-like comparison between a
manual polymorphic proof and a \emph{parametric} Trocq development, and to
demonstrate amortization to a second element type. On inspection its
parametric witnesses, though correct and machine-checked, are never applied:
the lemma intended to consume them is discharged by the manual polymorphic
proof instead. This is not an oversight but a type error in the original
design---the parametric witness concludes that summation agrees on two
\emph{related} lists, whereas the goal is an algebraic distributivity
statement whose right-hand side is not of that form. We therefore make no
claim about parametric amortization; the partial result in \texttt{bs\_a4.v},
which isolates the three strategies on \texttt{length}, is what the corpus
currently supports. Discharge requires restating the parametric route with a
goal the witnesses can prove, and re-measuring.

We consider O1--O5 open rather than fatal. Each bounds the scope of
Section~\ref{sec:results} to what the corpus exercises, and each has a
concrete stated experiment that would discharge it.


\section{Related Work}
\label{sec:related}

Trocq~\cite{cohen2024trocq,cohen2025trocq} and the univalent parametricity
framework it generalizes~\cite{tabareau2018equivalences,tabareau2021marriage}
are the object of this study; both focus on expressiveness and correctness,
and neither reports an effort comparison against a manual baseline, which is
the gap this paper addresses for one small corpus. CoqEAL's data-refinement
parametricity~\cite{cohen2013refinements} and Isabelle/HOL's Lifting and
Transfer~\cite{huffman2013lifting} solve closely related problems in
different settings; a natural extension is to check whether the regime
classification found here also organizes their cost trade-off. Proof
repair~\cite{ringer2018proof,ringer2021proofrepair}, and its extension to
quotient-type equivalences~\cite{viola2025proof}, targets a related economic
question---when it is cheaper to \emph{patch} a broken proof than to redo
it---and could be compared under the same methodology; automated repair of
verified libraries~\cite{gopinathan2023mostly} raises the same question at
library scale. On the empirical side, the QED-at-Large
survey~\cite{ringer2020qed} documents effort and maintenance across verified
systems without treating proof transfer as a costed decision. Finally, the
fixed/variable decomposition is inherited from software cost
estimation~\cite{boehm1981economics}; to our knowledge this paper is the
first to apply it to a proof-transfer tactic rather than to a software
project.


\section{Conclusion and Future Work}
\label{sec:conclusion}

Starting from a naive linear cost model that the corpus immediately
falsified, we arrived through three further refinements at a result we did
not expect. The question ``when does proof transfer pay off?'' has no single
answer, because the manual baseline it is measured against is not a single
thing. While two representations share an algebraic signature, transfer is
find-and-replace, costs almost nothing, and Trocq does not repay its setup
until roughly fifteen theorems. As soon as the signature diverges---a type
class, a differing law, an evaluation order---renaming stops producing
typechecking scripts, the manual baseline acquires a fixed cost of its own,
and Trocq repays within three or four. The engineering question is therefore
not ``how many theorems?'' but ``do the signatures agree?'', and that
question is answerable in minutes by attempting one rename.

Two findings run against our initial hypothesis and we report them as such:
complexity lowers the break-even but does not abolish it, and Trocq's setup
does not amortize across element types---a manual proof that can be stated
polymorphically generalizes for roughly one step per theorem, which no
relational infrastructure in our corpus matched.

The obligations in Section~\ref{sec:ttv} set the agenda: a larger,
multi-domain, multi-relatedness-class corpus contributed by more than one
proof engineer, so that $P_{\text{manual}}$, $W_{\text{trocq}}$ and $k$ can be
fitted statistically rather than by exhaustive measurement of a small corpus;
an independent effort proxy to discharge O1; and a restatement of the
parametric route to close O5. A further direction, motivated by the
software-product-line origin of the problem in
Section~\ref{sec:intro}, is to instantiate the model on a
feature-oriented/product-level pair of representations of a variability-aware
artifact, where both the shared vocabulary $f$ and the divergence count are
plausibly larger than in this study, and where the decision the model
supports---whether to invest in an automated transfer translation before
undertaking a large proof effort---has direct practical stakes.

\subsubsection*{Artifact availability.}
The Coq/Rocq sources, the counting script \texttt{\_count\_steps.py} and the
generated measurement tables are distributed with this paper; every constant
reported above can be regenerated by re-running the script over the corpus.

\bibliographystyle{splncs04}
\bibliography{references}

\end{document}
